\documentclass[10pt,a4paper]{amsart}
\usepackage[margin=19mm]{geometry}
\usepackage{amsmath,amssymb,mathtools}
\usepackage[scr=rsfs,cal=euler]{mathalfa}
\usepackage{xfrac}
\usepackage[unicode,hidelinks,bookmarks=false]{hyperref}
\newcommand{\ie}{i.e.\ }
\newcommand{\cf}{cf.\ }
\newcommand{\resp}{respectively\ }
\newcommand{\Subsection}{\subsection}
\providecommand{\nobreakdash}{\nobreak\hbox{-}}
\setlength{\emergencystretch}{2em}
\multlinegap0pt
\usepackage{bm}
\usepackage{bbm}
\usepackage{dsfont}
\usepackage{xspace}
\usepackage{cancel}
\usepackage{mathtools}
\usepackage{amsthm}
\makeatletter
\@ifundefined{theorem}{\newtheorem{theorem}{Theorem}}{}
\makeatother
\newcommand{\vertiii}[1]{{\left\vert\kern-0.25ex\left\vert\kern-0.25ex\left\vert #1 
		\right\vert\kern-0.25ex\right\vert\kern-0.25ex\right\vert}}
\title[Nitsche-based FEM for the Laplace eigenvalue problem: spectr]{Nitsche-based FEM for the Laplace eigenvalue problem: spectral approximation and a posteriori error analysis}
\author{Arbaz Khan, David Mora, Jesus Vellojin}
\date{Source: arXiv:2606.17052v1 (CC BY 4.0)}
\begin{document}
\maketitle

\noindent\emph{Source and licence.} Nitsche-based FEM for the Laplace eigenvalue problem: spectral approximation and a posteriori error analysis. Version arXiv:2606.17052v1 (CC BY 4.0), distributed under the Creative Commons Attribution 4.0 International licence, as recorded in the arXiv licence metadata for this exact version and shown on the abstract page https://arxiv.org/abs/2606.17052v1. Full rights notice: licenses/arxiv-2606.17052-CC-BY-4.0.txt.\par\smallskip

\noindent\textit{Editorial scope.} Counted unit: the Theorem whose source label is \texttt{None} in the article (source0.tex lines 978--986, source label \texttt{None}), delivered together with the article's own complete original proof of it (source0.tex lines 988--1086, one \texttt{proof} environment of 99 lines). No mathematics is written, repaired or re-derived: statement and proof are the article's own text line for line. Required context retained uncounted: satassum11 (lines 943--945); eq:lagrange-two-xi-level (lines 959--963). Every definition, display and label of those lines is retained unchanged. Omitted from this package and not redistributed: 1--942, 946--958, 964--977, 987--987, 1087--2065. Nothing inside a retained excerpt is omitted or altered: every retained body is delivered exactly as the article's lines are written, so no omission receipt of any kind accompanies this package. Citations: the retained excerpts carry no citation command at all, so no bibliography entry of the article is transcribed and no reference is redistributed. Reference adaptation: none was needed and none was made; every reference inside the delivered unit resolves inside it, so no unresolved reference or citation is produced. Printed slips kept literal: no printed word, symbol, hypothesis or number of the counted statement or of its proof was altered. Layout and class: the article's own class is \texttt{siamart250211} with packages including lipsum, amsmath, amssymb, amsfonts, latexsym, stmaryrd, graphicx, float, mathrsfs, color, epstopdf, algorithmic, multirow, diagbox; the wrapper is the lane's pinned \texttt{amsart} editorial wrapper at 10pt on A4, which restates the theorem-like environments the retained text needs and adds the article's own macro definitions that the retained text needs (\texttt{\textbackslash vertiii}), with extra wrapper packages (bm, bbm, dsfont, xspace, cancel, mathtools). The delivered numbering is the unit's own. Assets: no retained span contains a figure environment, a graphics-inclusion command, a PDF-page inclusion, a table or a TikZ picture; no image file is copied into this package, the package declares no asset dependency and no image file is redistributed with it. Licence: version arXiv:2606.17052v1 of this article is distributed under the Creative Commons Attribution 4.0 International licence, as recorded in the arXiv licence metadata for this exact version and shown on the abstract page https://arxiv.org/abs/2606.17052v1. A full rights notice is delivered at licenses/arxiv-2606.17052-CC-BY-4.0.txt. Characters: every retained excerpt is pure ASCII, so the delivered engine, XeLaTeX with its default Latin Modern text font, reports no missing character.
\begin{equation}\label{satassum11}
   \|u-u_{h/2}\|_{h/2}\le \gamma\|u-u_h\|_{h}.
\end{equation}

\begin{align}
&\sum_{K\in\mathcal T_{h/2}}h_K^{-2}\|\chi\|_{0,K}^2 + \sum_{F\in\mathcal E_{h/2}^0}h_F^{-1}\|\chi\|_{0,F}^2 +
\sum_{F\in\mathcal E_{h/2}} \left(h_F^{-1}\|\chi\|_{0,F}^2 + h_F\|\partial_n \chi\|_{0,F}^2 \right)\leq C \vertiii{v}_{h/2}^2.
\label{eq:lagrange-two-xi-level}
\end{align}

\begin{theorem}[Reliability]
Assume that the saturation property \eqref{satassum11} holds for the discrete
eigenfunction associated with the same simple eigenvalue branch. Then there exists a
constant $C_\alpha>0$, independent of $h$ but possibly depending on the fixed Nitsche
stabilization parameter $\alpha$, such that
\[
   \|u-u_h\|_h \leq C_\alpha\big(\eta_h(u_h)+\theta_h\big).
\]
\end{theorem}

\begin{proof}
By the triangle inequality and the saturation assumption \eqref{satassum11}, we have
\[
   \|u-u_h\|_h \leq \frac{1}{1-\gamma}\|u_{h/2}-u_h\|_{h/2}.
\]
Therefore, it is enough to bound $\|u_{h/2}-u_h\|_{h/2}$. Since the symmetric Nitsche
bilinear form is coercive on the refined mesh, there exists $v\in V_{h/2}$, with
$\|v\|_{h/2}=1$, such that
\[
   \|u_{h/2}-u_h\|_{h/2} \leq C_\alpha A_{h/2}(u_{h/2}-u_h,v).
\]
%Let $\overline{v}\in V_h$ be the Lagrange interpolant of $v$ on the coarse mesh and
%set $\chi:=v-\overline{v}$. 
Then
\begin{equation}\label{releq11}
A_{h/2}(u_{h/2}-u_h,v) =A_{h/2}(u_{h/2}-u_h,\chi) + A_{h/2}(u_{h/2}-u_h,\overline{v})=: W_1+W_2.
\end{equation}

We first estimate $W_1$. Since $u_{h/2}$ solves the discrete eigenvalue problem on
the refined mesh,
\[
   A_{h/2}(u_{h/2},\chi) = \lambda_{h/2}(u_{h/2},\chi)_{0,\Omega}.
\]
Thus,
\[
   W_1 = \lambda_{h/2}(u_{h/2},\chi)_{0,\Omega} - A_{h/2}(u_h,\chi).
\]
Adding and subtracting $\lambda_h(u_h,\chi)_{0,\Omega}$ gives
\[
   W_1 = (\lambda_{h/2}u_{h/2}-\lambda_hu_h,\chi)_{0,\Omega} + \left[\lambda_h(u_h,\chi)_{0,\Omega} -  A_{h/2}(u_h,\chi) \right].
\]
The first term is bounded by the higher-order contribution. Indeed, using the
definition of $\theta_h$, the Poincar\'e inequality and the interpolation estimate
\eqref{eq:lagrange-two-xi-level}, we get
\[
   |(\lambda_{h/2}u_{h/2}-\lambda_hu_h,\chi)_{0,\Omega}| \leq C\theta_h\|\chi\|_{0,\Omega} \leq  C\theta_h.
\]

For the second term, an elementwise integration by parts yields
\begin{align}
\lambda_h(u_h,\chi)_{0,\Omega}-A_{h/2}(u_h,\chi) =& \sum_{K\in\mathcal T_h}(R_K,\chi)_{0,K} -\sum_{F\in\mathcal E_h^0} \big([\![\nabla u_h\cdot n]\!],\chi\big)_{0,F} \nonumber\\
&+ \sum_{F\in\mathcal E_h} \Big( (\partial_n\chi,u_h)_{0,F} -\alpha h_F^{-1}(u_h,\chi)_{0,F} \Big),
\label{eq:residual-identity-reliability}
\end{align}
where $R_K:=\lambda_hu_h+\Delta u_h$. We now estimate each term in \eqref{eq:residual-identity-reliability}. By the Cauchy--Schwarz inequality and \eqref{eq:lagrange-two-xi-level},
\[
\left|\sum_{K\in\mathcal T_h}(R_K,\chi)_{0,K}\right| \leq \left( \sum_{K\in\mathcal T_h}h_K^2\|R_K\|_{0,K}^2 \right)^{1/2} \left(
\sum_{K\in\mathcal T_h}h_K^{-2}\|\chi\|_{0,K}^2 \right)^{1/2} \leq C\eta_h(u_h).\]
Similarly,
\[
\left| \sum_{F\in\mathcal E_h^0} \big([\![\nabla u_h\cdot n]\!],\chi\big)_{0,F} \right| \leq \left(\sum_{F\in\mathcal E_h^0}h_F
\|[\![\nabla u_h\cdot n]\!]\|_{0,F}^2 \right)^{1/2} \left( \sum_{F\in\mathcal E_h^0}h_F^{-1}\|\chi\|_{0,F}^2 \right)^{1/2} \leq C\eta_h(u_h).
\]
The boundary contribution is estimated by using the last two terms in
\eqref{eq:lagrange-two-xi-level}. Indeed,
\[
\begin{aligned}
&\left|\sum_{F\in\mathcal E_h} \Big( (\partial_n\chi,u_h)_{0,F} -\alpha h_F^{-1}(u_h,\chi)_{0,F} \Big) \right|\\
&\quad \le\left(\sum_{F\in\mathcal E_h}h_F^{-1}\|u_h\|_{0,F}^2\right)^{1/2}\left[\left(\sum_{F\in\mathcal E_h}h_F\|\partial_n\chi\|_{0,F}^2 \right)^{1/2} + \alpha \left( \sum_{F\in\mathcal E_h} h_F^{-1}\|\chi\|_{0,F}^2 \right)^{1/2} \right].
\end{aligned}
\]
By \eqref{eq:lagrange-two-xi-level} and the fact that
\(\|v\|_{h/2}=1\), the bracketed factor is bounded by a constant depending at most
on the fixed parameter \(\alpha\). Therefore,
\[
\left| \sum_{F\in\mathcal E_h} \Big( (\partial_n\chi,u_h)_{0,F} -\alpha h_F^{-1}(u_h,\chi)_{0,F} \Big) \right|\le C_\alpha \eta_\Gamma(u_h) \le C_\alpha \eta_h(u_h).
\]
Therefore,
\[
   |W_1| \leq C_\alpha\big(\eta_h(u_h)+\theta_h\big).
\]

We now estimate $W_2$. Since $\overline{v}\in V_h$, the discrete eigenvalue equation
on the coarse mesh gives
\[
   A_h(u_h,\overline{v}) = \lambda_h(u_h,\overline{v})_{0,\Omega}.
\]
Hence,
\begin{align*}
W_2 &= A_{h/2}(u_{h/2},\overline{v}) - A_{h/2}(u_h,\overline{v}) \\
&=\lambda_{h/2}(u_{h/2},\overline{v})_{0,\Omega} - \lambda_h(u_h,\overline{v})_{0,\Omega} + \big(A_h-A_{h/2}\big)(u_h,\overline{v}).
\end{align*}
The first two terms are bounded by $C\theta_h$, using the stability of the interpolant
and the normalization of $v$. The last term comes from the change of the Nitsche boundary
weights when passing from $\mathcal T_h$ to $\mathcal T_{h/2}$. Since $\alpha$ is fixed,
this contribution is bounded by
\[
   \left|\big(A_h-A_{h/2}\big)(u_h,\overline{v})\right| \leq  C_\alpha \eta_h(u_h).
\]
Consequently,
\[
   |W_2| \leq C_\alpha\big(\eta_h(u_h)+\theta_h\big).
\]
Combining the bounds for $W_1$ and $W_2$ in \eqref{releq11}, we obtain
\[
   \|u_{h/2}-u_h\|_{h/2} \leq C_\alpha\big(\eta_h(u_h)+\theta_h\big).
\]
The assertion follows from the saturation argument.
\end{proof}

\end{document}
